\documentclass[12pt,b5paper]{ltjsarticle}

%\usepackage[margin=15truemm, top=5truemm, bottom=5truemm]{geometry}
%\usepackage[margin=10truemm,left=15truemm]{geometry}
\usepackage[margin=10truemm]{geometry}

\usepackage{amsmath,amssymb}
%\pagestyle{headings}
\pagestyle{empty}

%\usepackage{listings,url}

% 記号変更
%\renewcommand{\theenumi}{(\arabic{enumi})}
\renewcommand{\labelenumi}{(\arabic{enumi})}


%\usepackage{graphicx}

%\usepackage{tikz} % draw graphics : TikZ ist kein Zeichenprogramm
%\usetikzlibrary{arrows.meta}

%\usepackage{wrapfig}
%\usepackage{bm} % ベクトルの矢印

%\usepackage{luatexja-ruby} % ルビを振る

%% 核Ker 像Im Hom を定義
\newcommand{\Ker}{\mathop{\mathrm{Ker}}\nolimits}
\newcommand{\Img}{\mathop{\mathrm{Img}}\nolimits}
%\newcommand{\Hom}{\mathop{\mathrm{Hom}}\nolimits}

%\DeclareMathOperator{\Rot}{rot}
%\DeclareMathOperator{\Div}{div}
%\DeclareMathOperator{\Grad}{grad}
%\DeclareMathOperator{\arcsinh}{arcsinh}
%\DeclareMathOperator{\arccosh}{arccosh}
%\DeclareMathOperator{\arctanh}{arctanh}

\usepackage{url} % URLの記述

%\usepackage{listings}
%
%\lstset{
%%プログラム言語(複数の言語に対応，C,C++も可)
%  language = Python,
%%  language = Lisp,
%%  language = C,
%  %背景色と透過度
%  %backgroundcolor={\color[gray]{.90}},
%  %枠外に行った時の自動改行
%  breaklines = true,
%  %自動改行後のインデント量(デフォルトでは20[pt])
%  breakindent = 10pt,
%  %標準の書体
%%  basicstyle = \ttfamily\scriptsize,
%  basicstyle = \ttfamily,
%  %コメントの書体
%%  commentstyle = {\itshape \color[cmyk]{1,0.4,1,0}},
%  %関数名等の色の設定
%  classoffset = 0,
%  %キーワード(int, ifなど)の書体
%%  keywordstyle = {\bfseries \color[cmyk]{0,1,0,0}},
%  %表示する文字の書体
%  %stringstyle = {\ttfamily \color[rgb]{0,0,1}},
%  %枠 "t"は上に線を記載, "T"は上に二重線を記載
%  %他オプション：leftline，topline，bottomline，lines，single，shadowbox
%  frame = TBrl,
%  %frameまでの間隔(行番号とプログラムの間)
%  framesep = 5pt,
%  %行番号の位置
%  numbers = left,
%  %行番号の間隔
%  stepnumber = 1,
%  %行番号の書体
%%  numberstyle = \tiny,
%  %タブの大きさ
%  tabsize = 4,
%  %キャプションの場所("tb"ならば上下両方に記載)
%  captionpos = t
%}

%\usepackage{cancel}
%\usepackage{bussproofs}
%\usepackage{proof}

\begin{document}


\hrulefill

$n\in \mathbb{N} =\{ n\in\mathbb{Z} \mid n>0 \}$
\begin{enumerate}
 \item
      恒等式$(n^{2}+1) - (n+2)(n-2) = 5$を利用し、
      $n+2$と$n^{2}+1$の公約数は
      1 または 5 に限ることを示せ。

      \dotfill

      $n+2$と$n^{2}+1$の公約数を$g$とする。
      このとき、
      $1\leq g \leq \min\{n+2,n^{2}+1\}$
      である。

      $m_{n+2},\;m_{n^{2}+1} \in \mathbb{N}$が存在し、
      $n+2=gm_{n+2},\;n^{2}+1=gm_{n^{2}+1}$
      と書ける。

      $(n^{2}+1) - (n+2)(n-2) = 5$より、
      \begin{align}
       n^{2}+1
       &= (n+2)(n-2) + 5\\
       &= gm_{n+2}(n-2) + 5
      \end{align}
      
      $g$は$n^{2}+1$の約数でもあるので、
      右辺$gm_{n+2}(n-2) + 5$の約数でもある。
      よって、$g$は$5$の約数でなければならず、
      $g=1$または$g=5$であることとなる。


      \hrulefill

 \item
      $n+2$と$n^{2}+1$が
      1以外に公約数を持つような自然数$n$をすべて求めよ。

      \dotfill

      $n+2$と$n^{2}+1$の公約数は$1$と$5$以外になく、
      1以外に公約数を持つということは
      最大公約数が 5 であるということである。
      そこで、$k\in\mathbb{N}$として、
      $n+2=5k$とする。
      これにより、
      $n=5k-2$である。

      これを$n^{2}+1$に代入する。
      \begin{equation}
       n^{2}+1
        = (5k-2)^{2}+1
        = 25k^{2} -20k +5
      \end{equation}

      よって、
      $n+2=5k \; (k\in\mathbb{N})$
      のとき、
      公約数5を持つ。

      \hrulefill

 \item
      $2n+1$と$n^{2}+1$が
      1以外に公約数を持つような自然数$n$をすべて求めよ。

      \dotfill

      $2n+1$は奇数であるので、$n^{2}+1$との公約数は奇数のみである。

      $n^{2}+1$は次のように変形できる。
      \begin{gather}
       n^{2}+1
        = \left( 2n+1 \right) \left( \frac{1}{2}n-\frac{1}{4} \right) +\frac{5}{4}\\
       4(n^{2}+1)
       = ( 2n+1 ) ( 2n-1 ) +5
      \end{gather}

      $2n+1$と$n^{2}+1$の公約数を$c$とする。
      $4(n^{2}+1)=( 2n+1 ) ( 2n-1 ) +5$より、
      $c$は右辺を割り切る数であり、
      $c$は$2n+1$を割り切るから
      $5$も割り切る数ということになる。

      よって、$c=1$または$c=5$ということになる。

      $2n+1 =5k \; (k\in\mathbb{N})$と置き、
      $n^{2}+1$を計算する。
      \begin{equation}
       n^{2}+1
        = \left( \frac{5k-1}{2} \right)^{2}+1
        = \frac{5}{4}(5k^{2}-2k+1)
      \end{equation}

      $5k^{2}-2k+1$が$4$で割り切れれば、
      $n^{2}+1$は$5$で割り切れる自然数となる。

      \begin{equation}
       5k^{2}-2k+1
        =4k^{2} + k^{2}-2k+1
        =4k^{2} + (k-1)^{2}
      \end{equation}

      上記式より、
      $5k^{2}-2k+1$が$4$で割り切れるためには
      $k-1$が偶数であればよい。

      $k-1=2k^{\prime}$とすれば、
      $2n+1=5(2k^{\prime}+1)$から
      $n=5k^{\prime} + 2$となる。
      $k\in\mathbb{N}$であるので、
      $k^{\prime}=0,1,2,\dots$である。
%      \begin{align}
%       n^{2}+1
%       &= \frac{5}{4}(5k^{2}-2k+1)
%       = \frac{5}{4}(4k^{2} + (k-1)^{2})\\
%       &= \frac{5}{4}(4(2k^{\prime}+1)^{2} + (2k^{\prime})^{2})
%       = 5(2k^{\prime}+1)^{2} + 5(k^{\prime})^{2}
%      \end{align}

      $2n+1,\; n^{2}+1$を確認すると
      次のように$5$の倍数であることがわかる。
      \begin{align}
       2n+1 &= 5(2k^{\prime} +1)\\
       n^{2}+1 &= (5k^{\prime}+2)^{2}+1
       = 5( 5(k^{\prime})^{2} +4k^{\prime} +1)
      \end{align}

      よって、
      $1$以外に公約数を持つのは次の数のときである。
      \begin{equation}
       n= 5m +2 \quad (m\in\mathbb{N}_{0}=\mathbb{N} \cup  \{0\})
      \end{equation}


      \hrulefill

\end{enumerate}


\hrulefill









\end{document}
